% =====================================================================
\chapter{Conclusiones}
\label{cap:conclusiones}

\section{Respuesta a los objetivos}

Se cumplieron cinco de los seis objetivos específicos (tabla~\ref{tab:objetivos}): el corpus, el
modelo de sentimiento, las variables de shock, el agrupamiento de las empresas —que mostró que no
hay familias de comportamiento más allá de la liquidez— y el modelado con desfases. No se
integraron los \textit{trending topics} de X/Twitter, y el DLM previsto quedó reemplazado por el
VAR, que ya incorpora los efectos rezagados.

\section{La hipótesis}

La hipótesis H1 sostenía que el shock informativo tiene un impacto medible y significativo
sobre el precio o el volumen de las acciones argentinas. \textbf{Con el tono de las noticias,
no se sostiene}: el tono no anticipó el retorno ni el volumen en el mercado, en las empresas,
en los sectores ni en ningún régimen, ni siquiera con la prueba más favorable, que usa solo las
notas publicadas antes de la apertura. Lo que apareció con claridad es la \textbf{relación
inversa}: el retorno del S\&P~Merval anticipa el tono de las notas siguientes, en 39 de las 47
empresas y en 15 de las 17 con tono propio, y el efecto se concentra en las notas publicadas la
misma noche, después del cierre. Los medios económicos argentinos, en el período estudiado,
contaron lo que el mercado ya había hecho.

Con la definición de shock del plan —muchas más notas que lo habitual— H1 encuentra un
\textbf{apoyo parcial, en el volumen}: cuando una empresa aparece en muchas más notas, ese día
su acción se opera más. Ese mayor volumen no viene con una dirección en el precio, no se
arrastra a los días siguientes y no se anticipa con las notas de la noche. La cantidad de
cobertura acompaña a la actividad del mercado, pero no la anticipa.

% ---------------------------------------------------------------------
\section{Los instrumentos}

Para llegar a esa respuesta hubo que construir cuatro instrumentos que no existían para el
mercado argentino, y cada uno es un aporte en sí mismo:

\begin{itemize}
  \item \textbf{El corpus de noticias:} 14.012 notas de siete portales, recolectadas con la
        Wayback Machine y los mapas del sitio, auditadas y depuradas, con cada nota asignada a
        la primera rueda en la que el mercado pudo reaccionar a ella.
  \item \textbf{FinBETO:} un termómetro del tono de la prensa económica argentina, con un F1
        macro de 0,685 en cuatro clases sobre noticias no vistas, 17,7 puntos por encima del
        modelo de referencia en tres clases, y un conjunto de oraciones etiquetadas por
        sentimiento. Combinado con el NER, permitió medir por primera vez el tono de cada
        empresa descontando el clima general.
  \item \textbf{El VIX argentino:} un índice diario de tensión del mercado, construido con
        volatilidad realizada, que identifica sin ambigüedad los shocks del período —la salida
        del cepo y las dos elecciones de 2025— y los ordena por tamaño.
  \item \textbf{Los regímenes de mercado:} una clasificación de cada rueda en Calmo, Moderado y
        Tenso, estable ante distintas semillas, que muestra que al estrés se entra de golpe y
        se sale de a poco.
\end{itemize}

El experimento que los usa es, a su vez, un \textit{pipeline} reproducible de punta a punta,
con controles automáticos, especificaciones escritas antes de correr los análisis y corrección
por pruebas múltiples. Que los instrumentos funcionen es lo que da peso al resultado: el
termómetro responde con claridad a la suba del mercado, de modo que si no anticipa los precios
no es porque no mida.

% ---------------------------------------------------------------------
\section{Limitaciones}

\begin{itemize}
  \item \textbf{Un año de datos y datos diarios.} Con 224 ruedas, un efecto chico de las
        noticias podría no detectarse; y los datos diarios no ordenan lo que pasa dentro de una
        rueda (algunas acciones cotizan también en Nueva York después del cierre local).
  \item \textbf{Medición del tono.} El F1 mide el acuerdo de FinBETO con las etiquetas de Gemma;
        el tono de una empresa es el de las oraciones que la mencionan, y muchas ruedas se apoyan
        en una o dos.
  \item \textbf{Corpus y precios.} La cobertura depende de qué capturó la Wayback Machine, el
        13,4\% de las notas tiene horas por defecto y los precios están en pesos nominales.
  \item \textbf{VIX y eventos.} El VIX argentino es simétrico y conserva un shock durante 20
        ruedas, y los estudios de eventos del mercado tienen entre 5 y 16 eventos.
\end{itemize}

% ---------------------------------------------------------------------
\section{Líneas futuras}

Integrar los \textit{trending topics} de X/Twitter, como preveía el plan; implementar el DLM como
complemento del VAR; repetir los modelos con precios intradiarios y en dólares contado con
liquidación; reemplazar la media de 20 ruedas del VIX por una media exponencial
\citep{riskmetrics1996}; ampliar el etiquetado con una muestra anotada por personas y probar MarIA
como modelo base \citep{gutierrez2022maria}; separar las noticias por tema; y extender la ventana
a varios años, para probar si el tono anticipa al mercado en las recesiones, como encontró
\citet{garcia2013}.
